\begin{remark}[The two conflicts, and the one optionality behind them]\label{rem:conflict-cause}
PromQL's table carries $2$ conflicts, both a shift/reduce on the left parenthesis, and
neither is an accident of our encoding. They sit in two states --- one reached after
\texttt{group\_left} and one after \texttt{group\_right} --- and have a single cause, so
an earlier version of this subsection reported one conflict rather than two. That was not
a reading of the table: Lark raises at the \emph{first} conflict in strict mode, so the
message the measurement returned said one however many there were, and the count is now
read off the resolution log instead (\texttt{parse.py}'s
\texttt{lalr\_conflict\_report}). The cause is that the vector-matching modifiers take an
\emph{optional} parenthesised label list, so after either of them an opening parenthesis
may begin that list or a parenthesised right operand:
\begin{center}\footnotesize
\texttt{a * on(x) group\_left (y) b}\qquad vs.\qquad\texttt{a * on(x) group\_left (b)}
\end{center}
One token of lookahead cannot distinguish them at that point, and the conflict comes
from the optionality rather than from our layering: making the label list mandatory
removes it entirely, which is how we identify the cause, but that also rejects
\texttt{group\_left b}, which the language allows. What this shows is that \emph{this}
grammar is not LALR(1), not that no LALR(1) grammar for PromQL exists --- the language
is not ambiguous, since the expression above has one parse, which Prometheus produces
and then rejects for want of an operand. We make no claim about the whole class of
grammars for this language; a formulation that defers the decision past the
parenthesis may well exist.

The resolution is to shift, reading the parenthesis as the label list. That rejects
\texttt{a * on(x) group\_left (b)}, even though a reader can see the intended parse, and
it is not a liberty we take: Prometheus~3.15.0 rejects the same expression by the same
route, and accepting it would make this grammar \emph{more} permissive than the language
it describes (Remark~\ref{rem:conflict-resolution}).
\end{remark}
